\section{2026-08-19 — Tres clases, y el corpus argentino como subproducto de la operación}

Al escribir las tres secciones que faltaban en el capítulo de metodología aparecieron dos contradicciones de fondo que venían arrastrándose desde abril y que obligaron a decidir antes de redactar.

La primera es el número de clases del clasificador. El capítulo de antecedentes ya tenía impreso un esquema de cuatro categorías: verdadero, falso, engañoso e indeterminado. Ese esquema no coincidía con la salida del primer módulo del sistema, que produce tres probabilidades, ni con ningún conjunto de datos disponible: la categoría \enquote{engañoso} no existe en LIAR, ni en FakeNewsNet, ni en el corpus en español, de modo que habría que haberla definido y anotado desde cero sin ejemplo previo alguno. Se adoptaron tres clases: verdadero, falso y sin verificar. La tercera no es un punto intermedio de una escala sino la ausencia de evidencia, y se modela como clase y no como umbral de confianza porque es la situación más frecuente en el contenido de circulación reciente, que es el objeto de interés del sistema.

De esa decisión se desprendió una aclaración que faltaba en todo el trabajo y que se incorporó al capítulo de antecedentes: las clases del clasificador no son los estados del veredicto que percibe el usuario. El veredicto lo produce el módulo de orquestación combinando la salida del clasificador con el contraste externo y con las señales de la cuenta, de modo que ambos conjuntos tienen cardinalidad y significado distintos. La confusión entre ambos explicaba buena parte de las inconsistencias detectadas.

La segunda contradicción es el origen del corpus argentino. El material de abril lo describía como una recolección independiente por extracción de contenido de medios y etiquetas provenientes de organizaciones de verificación. Ese diseño quedó inviable: la organización de verificación de referencia bloquea de forma deliberada a todo cliente que no sea un navegador, y eludir ese bloqueo trasladaría la discusión al artículo 153 bis del Código Penal. Se resolvió que el corpus se construya como subproducto de la operación del sistema, según establece el requerimiento que obliga a persistir cada publicación analizada junto con su resultado y su evidencia. La estructura de ese registro coincide con la de un ejemplo de entrenamiento.

Esta segunda decisión tiene una consecuencia que conviene dejar asentada, porque no es evidente: vuelve coherentes entre sí el capítulo de descripción y el de metodología. El argumento comercial sostiene que el activo diferencial se produce sin costo de adquisición de datos, porque son los usuarios quienes analizan voluntariamente. Un corpus construido por extracción de contenido de medios habría contradicho ese argumento en el mismo documento que lo enuncia.

Un tercer punto quedó registrado al escribir la sección de validación. Ninguno de los conjuntos de datos externos aporta ejemplos de la clase sin verificar, dado que todos ellos etiquetan afirmaciones que ya atravesaron un proceso de verificación. Es el sesgo de selección propio de los corpus construidos a partir de organizaciones de verificación. En consecuencia, esa clase solo puede aprenderse del corpus propio, circunstancia que convierte su construcción en una condición necesaria del diseño y no en una contribución accesoria, y que conviene enunciar así en la defensa.

\bigskip
